
\medskip

BAI est un triangle rectangle en A tel que BA $= 210$~cm et AI $= 155$~cm. 

\medskip

\begin{enumerate}
\item %Déterminer la longueur BI au cm près.
D'après le théorème de Pythagore si BAI est un triangle rectangle en A , alors BI$^2 = \text{BA}^2 + \text{AI}^2 = 210^2 + 155^2 = \np{44100} + \np{24025} = \np{68125}$.

On a donc BI $ = \sqrt{\np{68125}} \approx 261,008$, soit 261~(cm) à l'unité près.

\textbf{Rédiger la réponse en faisant apparaître les différentes étapes.}

L'immeuble de Joanne possède 15 vitres rectangulaires.

%Chaque vitre a pour longueur 210 cm et pour largeur 155 cm.
%
%Lors d'une préalerte cyclonique Joanne pose de l'adhésif sur les deux diagonales de chaque vitre de l'immeuble.
%\begin{center}
%\psset{unit=1cm,arrowsize=2pt 3}
%\begin{pspicture}(6,6.8)
%%\psgrid
%\psframe(1.5,1.8)(4.5,6.3)
%\psline(1.5,1.8)(4.5,6.3)\psline(4.5,1.8)(1.5,6.3)
%\uput[d](3,6.3){$155$~cm}\uput[r](4.5,4.05){210 cm}
%\rput(3,0.2){Une bande d'adhésif est assimilée à une diagonale du rectangle}
%\rput(3,6.6){Schéma de la situation}
%\psline{->}(3,0.6)(2.3,3)\psline{->}(3,0.6)(3.7,3)
%\end{pspicture}
%\end{center}
%
\item %Justifier que Joanne a besoin d'environ $5,22$ m d'adhésif pour une vitre.
Sur chacune des 15 vitres Joanne doit poser deux bandes de 261 cm : elle pose donc sur chaque vitre  :

$2 \times 261 = 522$~(cm) d'adhésif, soit 5,22~(m).
\end{enumerate}

%Joanne a 7 rouleaux d'adhésif de $10$ m chacun.

\begin{enumerate}[resume]
\item %A-t-elle assez d'adhésif pour toutes les vitres ? Justifier la réponse.
Elle doit poser sur la totalité des 15 fenêtres :

$15 \times 5,22 = 78,3$~(m), alors qu'elle n'a que $7 \times 10 = 70$~(m). Il lui manque donc un rouleau.
\end{enumerate}

\bigskip

